% ECONOMICS_PROBLEM: {"id": "G12-455", "title": "Asset-demand elasticity", "jel_code": "G12", "source_id": "OP-0312"}
% !TeX program = xelatex
\documentclass[11pt,a4paper]{article}
\usepackage[margin=23mm]{geometry}
\usepackage{fontspec}
\setmainfont{Latin Modern Roman}
\usepackage{amsmath,amssymb,mathtools}
\usepackage{xcolor,enumitem,booktabs,longtable,array,xurl,hyperref}
\definecolor{muted}{HTML}{53636C}
\hypersetup{colorlinks=true,urlcolor=blue,linkcolor=blue}
\setlength{\parindent}{0pt}
\setlength{\parskip}{5pt}
\newcommand{\R}{\mathbb R}
\newcommand{\E}{\mathbb E}
\newcommand{\N}{\mathbb N}
\newcommand{\ind}{\mathbf 1}
\newcommand{\Var}{\operatorname{Var}}
\newcommand{\Cov}{\operatorname{Cov}}
\newcommand{\supp}{\operatorname{supp}}
\DeclareMathOperator*{\argmin}{arg\,min}
\DeclareMathOperator*{\argmax}{arg\,max}
\newcommand{\entrymeta}[3]{{\sffamily\small\color{muted}\textbf{\texttt{#1}}\quad|\quad #2\par #3\par}}
\newcommand{\Problemref}[1]{\texttt{#1}}
\newcommand{\JELref}[2]{\texttt{#2} (JEL \texttt{#1})}
\begin{document}
\section*{Asset-demand elasticity}
\textbf{Problem ID:} \texttt{G12-455}\quad \textbf{JEL:} \texttt{G12}\par
% BEGIN ECONOMICS STATEMENT
\label{problem:OP-0312}
{\sffamily\small\color{muted}Original subject group: Asset pricing and financial markets\par}
\label{definitions:problem:30007594}
\textbf{Audit classification:} Explicit published open question. The primary revision explicitly requests better identification of aggregate elasticity. The catalogue intervention and reciprocal residual-demand normalization are editorial, and the linked draft is a 2022 revision of the 2021 WP.
\subsection*{Definitions and precise research target}
Fix an asset market, horizon \(h\geq0\), publicly observed state \(x\), and pre-intervention fundamental information \(\mathcal F_t\). Let \(u\) be an externally induced net demand shift measured as a fraction of the market's initial value. Define \(p_{t+h}(u)\) as the logarithm of the equilibrium price under that intervention, including other investors' endogenous responses. The causal price multiplier is
\[\mu_h(x)=\left.\frac{\partial E[p_{t+h}(u)-p_{t+h}(0)\mid x]}{\partial u}\right|_{u=0}.\]
Where \(\mu_0(x)>0\), define the corresponding local residual-demand elasticity as \(\zeta(x)=1/\mu_0(x)\), under the stated normalization and fixed supply. The research task is to identify \(\mu_h(x)\) from portfolio holdings, flows and prices while separating demand shifts from news about cash flows and discount rates. An admissible instrument must shift demand, have no direct fundamental-price effect conditional on \(\mathcal F_t\), and account for correlated flows across investors. Establish when these conditions are credible, derive uncertainty or bounds when exclusion fails, and test variation across stress states, asset classes and horizons. Cross-sectional relative elasticities must not be equated with aggregate elasticities without additional restrictions. A single fitted multiplier does not answer the full state-dependent identification question.

\textbf{Additional formal definitions and scope (editorial).} Write \(Q>0\) for fixed share supply and \(q(p)\) for other investors' residual share demand as a differentiable function of log price. Normalize the intervention as an additional demand of \(uQ\) shares, or an equivalent initial-value dollar injection with the same derivative at zero. Clearing is \(q(p(u))+uQ=Q\). At \(q(p(0))=Q\), define \(\zeta=-q'(p(0))/Q>0\); the implicit function theorem gives \(p'(0)=1/\zeta\). This reciprocal relation requires this one-dimensional residual-demand normalization and does not apply to an arbitrary flow multiplier. Specify \(h\in\mathbb N_0\), the initial state \(x\), and conditional differentiability and domination allowing differentiation inside expectations. An instrument \(Z\) requires nonzero conditional covariance with the induced flow, a stated exclusion restriction on the structural price equation, and a declared independence or conditional-moment condition; relevance and exclusion are distinct assumptions.

For the identification part, declare an admissible class \(\Theta\) of structural models, including preferences, technologies, shock laws, measurement/sampling rules and any equilibrium selector. If \(O\) denotes the complete observable history and \(T(\theta)\) the target defined above, the identified set is \(\mathcal I_T=\{T(\theta):\theta\in\Theta,\ \mathcal L_\theta(O)=\mathcal L_{\rm obs}(O)\}\); point identification means this set is a singleton. A \(do(a)\) comparison replaces stated policy/technology equations while fixing the declared exogenous law and re-solving the remaining equations. These definitions do not assert that an identification theorem holds without further restrictions.
\subsection*{Variants and assumption changes}
\begin{enumerate}
\item \textbf{Single asset, fixed supply (editorial).} Identify \(\zeta\) in the scalar clearing equation from an excluded demand shifter with nonzero first stage.
\item \textbf{Aggregate and dynamic demand (source-motivated extension).} Allow cross-asset substitution and adjustment delays; identify the full price-response matrix and horizon-dependent responses, whose aggregate inverse need not equal cross-sectional elasticities.
\end{enumerate}
\subsection*{References and source audit}
\begin{enumerate}
\item NBER record verifies 2021 WP 28967 and DOI. Linked Cowles PDF cover dates this revision May 12, 2022. Locator: Section 7, printed p.44 (PDF p.44), last paragraph on alternative identification strategies. Support: Elasticity identification agenda. Full 2022 revision accessible. URL: \url{https://cowles.yale.edu/sites/default/files/2022-10/SSRN-id3686935.pdf}.
\item Publisher metadata and abstract verified. Locator: abstract. Background only; no exact open passage verified. URL: \url{https://www.annualreviews.org/content/journals/10.1146/annurev-financial-090524-120754}.
\end{enumerate}
\begin{enumerate}\raggedright
\item\hypertarget{definitions:source:30007594:1}{} Xavier Gabaix; Ralph S. J. Koijen. 2021. \emph{In Search of the Origins of Financial Fluctuations: The Inelastic Markets Hypothesis}. Section 7, printed p.44 (PDF p.44), last paragraph on alternative identification strategies.
\url{https://cowles.yale.edu/sites/default/files/2022-10/SSRN-id3686935.pdf}.
DOI: \url{https://doi.org/10.3386/w28967}.
\item\hypertarget{definitions:source:30007594:2}{} Valentin Haddad; Tyler Muir. 2025. \emph{Market Macrostructure: Institutions and Asset Prices}. Publisher abstract.
\url{https://www.annualreviews.org/content/journals/10.1146/annurev-financial-090524-120754}.
DOI: \url{https://doi.org/10.1146/annurev-financial-090524-120754}.
\end{enumerate}

\subsection*{Common conventions: Source mathematical conventions}

These conventions apply to each entry unless explicitly replaced.

\textbf{Models and identification.} A primitive model \(m\) specifies all probability laws, feasible choices, information, equilibrium and measurement rules used by its target. The admissible class \(\mathcal M\) is fixed independently of outcomes. If \(P_O(m)\) is its observable probability law and \(\tau(m)\) the target, its identified set at observable law \(P\) is
\[
 \mathcal I_\tau(P)=\{\tau(m):m\in\mathcal M,\ P_O(m)=P\}.
\]
Point identification means that this set is a singleton. An empty set signals incompatibility of data and assumptions. Coordinate bounds need not describe the joint set. Bounds are sharp only if every retained target is attainable or arbitrarily approachable by compatible models. Finite bounds require explicit restrictions on unobserved quantities; unrestricted confounding, hidden wealth, extrapolation or arbitrary equilibrium selection can yield uninformative sets.

\textbf{Causal interventions.} A potential outcome \(Y_i(p)\) is the outcome under a complete policy regime \(p\), including timing, financing, information and market spillovers. The target population and units are fixed before treatment. With interference, use the complete assignment vector or a justified exposure mapping; individual no-interference notation alone cannot identify an economy-wide intervention. A difference of mean potential outcomes does not require the cross-world joint distribution. Conditioning on post-treatment migration, survival, enrollment or employment changes the estimand and needs separate justification.

\textbf{Statistical statements.} An estimator, confidence set, forecast or test specifies a sampling law, model class and loss/coverage criterion. Uniform \((1-\alpha)\)-coverage means \(\inf_{m\in\mathcal M}\Pr_m\{\tau(m)\in C_n\}\geq1-\alpha\), or an explicitly asymptotic version. The whole parameter space trivially has coverage, so informativeness requires a stated width, risk or power objective. A forecasting improvement is relative to a named benchmark, information set and evaluation population.

\textbf{Equilibrium and optimization.} An equilibrium comprises feasible plans, optimality, beliefs where needed, budget constraints and all market/resource clearing. The primitives or a stated benchmark define each correspondence \(\mathcal E(p;m)\). Nonemptiness is assumed or proved, never inferred just from compact policy sets. A selected-equilibrium target uses a declared selection rule; a robust target quantifies over all permitted equilibria. Use suprema or infima unless attainment is established. An \(\varepsilon\)-optimal policy is feasible and within \(\varepsilon\) of the finite optimum value. Report an empty feasible set separately.

\textbf{Welfare and accounting.} Normative utility, interpersonal normalization, social weights, numeraire, discounting and the use of public receipts are primitives. Behavioral utility need not coincide with the welfare criterion. Household welfare includes ownership income and rents once; adding profits again to compensating variation generally double-counts them. A monetary equivalent is defined by an explicit transfer and price convention, with monotonicity and range conditions. Average effects, mechanism contributions, transfers and welfare losses are different targets. Infinite sums require convergence and dynamic feasible plans require terminal or no-Ponzi conditions.

\textbf{Source versions.} A working-paper date, revision date and journal publication year are distinguished. A DOI for a journal version is not automatically the DOI of an earlier manuscript. Locators refer to the stated version and printed pages where specified. A metadata-only check does not verify a claim about a full-text conclusion. Every entry's source subsection narrows the provenance claim accordingly.
% END ECONOMICS STATEMENT
\end{document}
